%============================================================
%    2 D   N U M E R I C S
%
\section{Numerical test cases}  \label{sec:numerics}
In this paper only numerical tests related to the two-dimensional Euler equations are considered. Obviously, CATMOOD schemes could be applied to any systems of conservation laws without any restriction.
Our methodology of testing relies on several classical and demanding test cases: 
\begin{enumerate}
    \item \textit{Isentropic vortex in motion}. This test measures the ability of the MOOD procedure combined with CAT schemes to achieve the optimal high order for a smooth solution. We also compared CATMOOD6 with unlimited CAT schemes, limited ACAT ones and some first order scheme using Rusanov, HLL and HLLC fluxes.
    \item \textit{Sedov Blast wave.} This problem has an exact solution presenting a single cylindrical shock wave followed by an exponential decay \cite{Sedov}. This test is used to check the behaviour of the CATMOOD6 scheme against shocks and to compare CATMOOD6 versus ACAT6.
    \item \textit{2D Riemann problems}. We simulate four versions of the four-state 2D Riemann problems \cite{RP_Schultz93}. These problems present large smooth regions, unstable shear layers, unaligned contact discontinuities and shock waves, along with complex interaction patches for which no exact solution has yet been derived.
    \item \textit{Astrophysical jet}. This intense and demanding test case challenges the positivity and robustness of the CATMOOD6 scheme as it presents an extremely violent jet at Mach 2000 generating a bow shock and unstable shear layers. 
\end{enumerate}
\begin{comment}
\rred{The usual CFL condition in 2D seems to be
$$
\Delta t^n \left( \frac{\Delta x}{\lambda_{ {\rm max}_x}^n} + \frac{\Delta y}{\lambda_{ {\rm max}_y}^n}\right) = \mathrm{CFL}\,  $$
or the more restrictive one
$$
\Delta t^n = \frac{\mathrm{CFL}}{2}\, \min\left( \frac{\Delta x}{\lambda_{ {\rm max}_x}^n}, \frac{\Delta y}{\lambda_{ {\rm max}_y}^n}\right), $$
what implies that CFL = 0.4 with the formula below is really CFL = 0.8.}
\end{comment}
The time step is chosen according to 
$$
\Delta t^n = \mathrm{CFL}\, \min\left( \frac{\Delta x}{\lambda_{ {\rm max}_x}^n}, \frac{\Delta y}{\lambda_{ {\rm max}_y}^n}\right),   $$
where $\lambda_{{\rm max}_x}^n$ and $\lambda_{{\rm max}_y}^n$ represent the maximum, over cells, of the spectral radius of the Jacobian matrices $\partial \mathbf{F}/\partial u$ and $\partial \mathbf{G}/\partial u$. 
Empirically we found that the CFL number should be less or equal to $0.5$.
In our calculations we choose CFL$=0.4$.
In all the numerical examples below (except the one in Section \ref{ssec:Mach}), the units are chosen in such a way that the order of magnitude of all field quantities is one, therefore we always use the same value for the constants appearing in \eqref{eq:delta}, i.e.\ $\varepsilon_1=10^{-4}$ and $\varepsilon_2=10^{-3}$, while the value $\varepsilon = 10^{-8}$ has been used in \eqref{wavefuction}.


 

%---------------------------------------------
%  2 D   V O R T E X     T E S T     C A S E   
\subsection{Isentropic vortex in motion} \label{ssec:vortex}
The isentropic vortex problem \cite{Shu1} challenges the accuracy of numerical methods since an exact, smooth and analytic solution exists. The computational domain is set to $\Omega = [-10,10]\times [-10,10]$.
The ambient flow is characterized by $\rho_\infty = 1.0$, $u_\infty = 1.0$, $v_\infty = 1.0$ and $p_\infty = 1.0$,
with a normalized ambient temperature $T^*_\infty =1.0$. 
At the initial time, $t=0$, onto this ambient flow is superimposed a vortex centered at $(0,0)$ with the following state:
$u = u_\infty + \delta u$, $ v = v_\infty + \delta v$, 
$ T^* = T^*_\infty + \delta T^*$, where the increments are given by
\begin{equation}
    \begin{array}{c}
      \nonumber
      \delta u   = -y' {\dfrac {\beta} {2 \pi}} \exp \left( {\dfrac {1-r^2} {2}} \right),  \quad 
      \delta v   = x' {\dfrac {\beta} {2 \pi}} \exp \left( {\dfrac {1-r^2} {2}} \right), \quad
      \delta T = - { \dfrac {(\gamma - 1 ) \beta^2} {8 \gamma \pi^2}} \exp \left( {1-r^2} \right),
    \end{array}
\end{equation}
with $r = \sqrt{x^2 + y^2}$. The so-called strength of the vortex is set to $\beta=5.0$ and the initial density is
given by  $\rho =   \rho_\infty \left( {{T}/{T_\infty} } \right)^{\frac{1}{\gamma-1} }$. Periodic boundary conditions are prescribed. 
At final time $t=t_{\text{final}}=20$ the vortex is back to its original position, and, the final exact solution matches the initial one. 
% What we expect.
Since the solution is smooth, it should be simulated with optimal high accuracy, in other words, the limiting/stabilization procedure employed in the scheme should not have any effect. 
% --- TAB ---
\begin{table}[!ht]
\numerikNine
    \centering
    \begin{tabular}{|c||cc|cc|cc|cc|}
    \hline \multicolumn{9}{|c|}{\textbf{2D Isentropic Vortex in motion - Rate of convergence}} \\
\hline   
\hline   
& \multicolumn{2}{c|}{\textbf{Rusanov-flux}} & \multicolumn{2}{c|}{\textbf{HLL}} & \multicolumn{2}{c|}{\textbf{HLLC}} & \multicolumn{2}{c|}{\textbf{CATMOOD6}}\\
 $N$ &  $L^1$ error &  order & $L^1$ error &  order & $L^1$ error &  order & $L^1$ error &  order \\ \hline
50 $\times$ 50  & 8.44$\times 10^{-3}$  &  ---    & 8.44$\times 10^{-3}$ & ---   &7.91$\times 10^{-3}$  &  -& 8.48$\times 10^{-3}$ & ---     \\
100 $\times$ 100 & 8.04$\times 10^{-3}$ &  0.07 & 8.04$\times 10^{-3}$ & 0.07&6.86$\times 10^{-3}$ &  0.21& 3.77$\times 10^{-3}$ & 1.17 \\
200 $\times$ 200 & 6.68$\times 10^{-3}$  & 0.27 & 6.67$\times 10^{-3}$ & 0.27&5.31$\times 10^{-3}$ &  0.37& 2.40$\times 10^{-7}$ & 13.94  \\
300 $\times$ 300 & 5.71$\times 10^{-3}$  & 0.36 & 5.71$\times 10^{-3}$ & 0.36&4.53$\times 10^{-3}$ &  0.39& 2.06$\times 10^{-8}$ & 6.05   \\
400 $\times$ 400 & 4.98$\times 10^{-3}$ &  0.47 & 4.98$\times 10^{-3}$ & 0.47&3.86$\times 10^{-3}$ &  0.55& 3.52$\times 10^{-9}$ & 6.14   \\
& \text{Expected} & 1 
& \text{Expected} & 1 
& \text{Expected} & 1 
& \text{Expected} & 6
\\
\hline  \hline  
& \multicolumn{2}{c|}{\textbf{CAT2}} & \multicolumn{2}{c|}{\textbf{CAT4}} & \multicolumn{2}{c|}{\textbf{CAT6}} &\multicolumn{2}{c|}{\textbf{ACAT6}} \\ 
     $N$ &  $L^1$ error &  order & $L^1$ error &  order & $L^1$ error &  order & $L^1$ error &  order  \\ \hline
50 $\times$ 50  & 7.94$\times 10^{-3}$  &  ---    & 2.03$\times 10^{-3}$ & ---   &8.46$\times 10^{-4}$  &  -& 8.95$\times 10^{-3}$ & --- \\
100 $\times$ 100 & 2.55$\times 10^{-3}$ &  1.64 & 1.42$\times 10^{-4}$ & 3.83&1.56$\times 10^{-5}$ &  5.76& 8.28$\times 10^{-3}$ & 0.11\\
200 $\times$ 200 & 6.12$\times 10^{-4}$ &  2.06 & 8.34$\times 10^{-6}$ & 4.09&2.41$\times 10^{-7}$ &  6.02& 8.34$\times 10^{-5}$ & 9.95\\
300 $\times$ 300 & 2.69$\times 10^{-4}$ &  2.02 & 1.64$\times 10^{-6}$ & 4.02&2.09$\times 10^{-8}$ &  6.03& 1.05$\times 10^{-5}$ & 5.14\\
400 $\times$ 400 & 1.52$\times 10^{-4}$ &  1.99 & 5.16$\times 10^{-7}$ & 4.01&3.68$\times 10^{-9}$ &  6.03& 2.48$\times 10^{-6}$ & 4.93\\
& \text{Expected} & 2
& \text{Expected} & 4 
& \text{Expected} & 6 
& \text{Expected} & 6
\\
    \hline 
    \end{tabular}
    \caption{Isentropic vortex in motion  $L^1-$norm errors on density $\rho$ between the numerical solution and the exact solution of the isentropic vortex in motion problem at $t_{\text{final}} = 20$ on uniform Cartesian mesh. }
    \label{tab:vortex_convergence}
\end{table}
% --- TAB ---

% --- TAB ---
\begin{table}[!ht]
    \centering
    \numerikNine
 \resizebox{\columnwidth}{!}{
    \begin{tabular}{|c||c| c| c| c|c|c|c|c|}
    \hline 
    & \multicolumn{3}{c|}{1st order methods} & 2nd ord. & 4th ord. & \multicolumn{3}{c|}{6th order methods} \\
    \cline{2-9}
    & \textbf{Rusanov} & \textbf{HLL} &  \textbf{HLLC}   &  \textbf{CAT2} & 
    \textbf{CAT4} & \textbf{CAT6} &  \textbf{CATMOOD6}   &  \textbf{ACAT6} \\
    \hline \hline
    \textbf{CPU(s)} & 48.22 & 60.58   & 98.41 & 424.72  & 4375.68 & 16112.38   & 20848.26 & 23597.67    \\ \hline
    \textbf{Ratio} & 1 & 1.26  & 2.04 &  8.81 & 90.74 & 334.14  & 432.36 &  489.38\\ 
    \hline
    \end{tabular} 
    }
    \caption{Isentropic vortex in motion with $300\times300$ cells. First row: computational costs expressed in seconds. Second row: ratio with respect to the Rusanov cost.}
    \label{tab:vortex_cpu}
\end{table}
% --- TAB ---
\begin{comment}
% --- TAB ---
\begin{table}[!ht]
    \centering
    \begin{tabular}{|c| c| c| c|}
    \hline 
    \rm{Rusanov flux} & \rm{HLL} &  \rm{HLLC}   &  \rm{CAT2} \\
    \hline 
    48.22 & 60.58   & 98.41 & 424.72   \\ \hline
    1 & 1.26  & 2.04 &  8.81 \\ \hline
    \hline
    \rm{CAT4} & \rm{CAT6} &  \rm{CATMOOD6}   &  \rm{ACAT6} \\
    \hline 
    4375.68 & 16112.38   & 20848.26 & 23597.67   \\ \hline
    90.74 & 334.14  & 432.36 &  489.38 \\
    \hline
    \end{tabular} 
    \caption{Isentropic vortex in motion. First row: computational costs express in second. Second row: ratio between the computational costs of the methods and the Rusanov scheme.}
    \label{tab:vortex_cpu}
\end{table}
% --- TAB ---
\end{comment}
\begin{comment}
% --- TAB ---
\begin{table}[!ht]
\numerikNine
    \centering
    \begin{tabular}{|c|c|c|c|c|}
    \hline \multicolumn{5}{|c|}{\textbf{2D Isentropic Vortex in motion - Rate of convergence}} \\
    \hline   & \multicolumn{2}{c|}{\textbf{Rusanov-flux}} & \multicolumn{2}{c|}{\textbf{HLL}} \\
    \hline  $N$ &  $L^1$ error &  order & $L^1$ error &  order \\ \hline
            $50 \times 50$  & $8.44E-03$  &  $-   $ & $8.44E-03$ & $-$  \\
            $100 \times 100$ & $8.04E-03$ &  $0.07$ & $8.04E-03$ & $0.07$ \\
            $200 \times 200$ & $6.68E-03$  & $0.27$ & $6.67E-03$ & $0.27$  \\
            $300 \times 300$ & $5.71E-03$  & $0.36$ & $5.71E-03$ & $0.36$  \\
            $400 \times 400$ & $4.98E-03$ &  $0.47$ & $4.98E-03$ & $0.47$   \\
    \hline  & \multicolumn{2}{c|}{\textbf{HLLC}} & \multicolumn{2}{c|}{\textbf{CATMOOD6}} \\ \hline
     $N$ &  $L^1$ error &  order & $L^1$ error &  order \\ \hline   
            $50 \times 50$  & $7.91E-03$  &  $-$    & $8.48E-03$ & - \\
            $100 \times 100$ & $6.86E-03$ &  $0.21$ & $3.77E-03$ & $1.17$ \\
            $200 \times 200$ & $5.31E-03$ &  $0.37$ & $2.40E-07$ & $13.94$ \\
            $300 \times 300$ & $4.53E-03$ &  $0.39$ & $2.06E-08$ & $6.05$ \\
            $400 \times 400$ & $3.86E-03$ &  $0.55$ & $3.52E-09$ & $6.14$ \\
     \hline  & \multicolumn{2}{c|}{\textbf{CAT2}} & \multicolumn{2}{c|}{\textbf{CAT4}} \\ \hline
     $N$ &  $L^1$ error &  order & $L^1$ error &  order \\ \hline   
            $50 \times 50$  & $7.94E-03$  &  $-$    & $2.03E-03$ & $-$ \\
            $100 \times 100$ & $2.55E-03$ &  $1.64$ & $1.42E-04$ & $3.83$ \\
            $200 \times 200$ & $6.12E-04$ &  $2.06$ & $8.34E-06$ & $4.09$ \\
            $300 \times 300$ & $2.69E-04$ &  $2.02$ & $1.64E-06$ & $4.02$ \\
            $400 \times 400$ & $1.52E-04$ &  $1.99$ & $5.16E-07$ & $4.01$ \\
     \hline  & \multicolumn{2}{c|}{\textbf{CAT6}} & \multicolumn{2}{c|}{\textbf{ACAT6}} \\ \hline
     $N$ &  $L^1$ error &  order & $L^1$ error &  order \\ \hline   
            $50 \times 50$  & $8.46E-04$  &  $-$    & $8.95E-03$ & $-$ \\
            $100 \times 100$ & $1.56E-05$ &  $5.76$ & $8.28E-03$ & $0.11$ \\
            $200 \times 200$ & $2.41E-07$ &  $6.02$ & $8.34E-05$ & $9.95$ \\
            $300 \times 300$ & $2.09E-08$ &  $6.03$ & $1.05E-05$ & $5.14$ \\
            $400 \times 400$ & $3.68E-09$ &  $6.03$ & $2.48E-06$ & $4.93$ \\
    \hline 
    \end{tabular}
    \caption{Isentropic vortex in motion  $L^1-$norm errors on density $\rho$ between the numerical solution and the exact solution of the isentropic vortex in motion problem at $t_{final} = 20$ on uniform Cartesian mesh. }
    \label{tab:vortex_convergence}
\end{table}
\end{comment}
% --- FIG ---
\begin{figure}[!ht]
    \centering    
    \includegraphics[width=0.495\textwidth]{FIG/Num_Exp/Vortex/Vortex_CM6.eps}
    \includegraphics[width=0.495\textwidth]{FIG/Num_Exp/Vortex/errors.pdf}
    \caption{2D isentropic vortex in motion ---
    Left: Numerical solution at final time for density with CATMOOD6 with HLLC flux scheme used for the first order method on $100 \times 100$ uniform mesh and CFL$=0.4$ ---
    Right: Errors in $L_1$ norms given by all schemes.}
   \label{fig:vortex_density}
\end{figure}
% --- FIG ---
We run the isentropic vortex test case with first order Rusanov-flux, HLL, HLLC and  sixth order CATMOOD method on successive refined Cartesian meshes going from $50 \times 50$  up to  $400 \times 400$ cells. 
The density at final time is plotted in figure~\ref{fig:vortex_density}-left.
The results of the convergence analysis for all the schemes are displayed in table~\ref{tab:vortex_convergence} and in figure~\ref{fig:vortex_density}-right.
The expected rate of convergence is reached for CAT schemes, while for the first-order schemes the convergence is below the expected order.
Ideally the limited $6$th-order ACAT6 scheme should produce the same errors as the ones given by CAT6 scheme, but ACAT6 errors are two orders of magnitude greater.
Contrarily, starting at mesh $200\times 200$, CATMOOD6 errors match the optimal ones from CAT6, and, the nominal $6$th order is retrieved. This proves that the \aposteriori MOOD limiting avoids spurious intervention if the mesh is fine enough.
%\ema{
Table~\ref{tab:vortex_cpu} presents the computational costs of various methods for  the isentropic vortex test case on a uniform grid with $300\times300$ cells, CFL$ = 0.4,$ final time $t_{\text{final}} = 20$, and periodic boundary conditions. The first row displays the CPU computational costs in seconds for the eight methods, while the second row shows the ratio between the computational cost of each scheme and the one of the Rusanov method. This table emphasizes the superiority of the \aposteriori approach over the \apriori one. Indeed ACAT6 scheme appears to be not only less accurate than CATMOOD6, but also its execution time is higher.
%}


%---------------------------------------------
%    2D   S E D O V    B L A S T W A V E  

\subsection{Sedov blast wave}\label{ssec:Sedov}
%--- FIG ---
\begin{figure}[!ht]
     \centering
     \begin{subfigure}[b]{0.49\textwidth}
         \centering
         \includegraphics[width=\textwidth]{FIG/Num_Exp/Sedov/CATMOOD.eps}
         \caption{Projection of the numerical density for CATMOOD6 with HLLC as first order scheme}
         \label{Sed:CATMOOD6}
     \end{subfigure}
     \hfill
     \begin{subfigure}[b]{0.49\textwidth}
         \centering
         \includegraphics[width=\textwidth]{FIG/Num_Exp/Sedov/CATMOOD_cut.eps}
         \vspace{0.2cm}
         \caption{Zoom on $[0,1.2]\times[0,1.2]$ of the numerical density for CATMOOD6.}
         \label{Sed:CATMOOD6_cut}
     \end{subfigure}
     \caption{Sedov blast wave \ref{ssec:Sedov}. Numerical density obtained with CATMOOD6 with HLLC as first order scheme on the interval $[-1.2,1.2]\times[-1.2,1.2]$ adopting a $200\times200$ mesh and CFL$=0.4$ at time $t=1$. Projection on the plane O-$xy$ (left); zoom on the interval $[0,1.2]\times[0,1.2]$ (right).}
     \label{Se:CATMOOD6}
\end{figure}

\begin{comment}
\begin{figure}[!ht]
     \centering
     \begin{subfigure}[b]{0.4\textwidth}
         \centering
         \includegraphics[width=\textwidth]{FIG/Num_Exp/Sedov/HLL_100.eps}
         \caption{1st order scheme on a $100\times100$ mesh.}
         \label{Sed:HLLC_100}
     \end{subfigure}
     \hfill
     \begin{subfigure}[b]{0.4\textwidth}
         \centering
         \includegraphics[width=\textwidth]{FIG/Num_Exp/Sedov/HLL_200.eps}
         \caption{1st order scheme on a $200\times200$ mesh.}
         \label{Sed:HLLC}
     \end{subfigure}
     \\
     \begin{subfigure}[b]{0.4\textwidth}
         \centering
         \includegraphics[width=\textwidth]{FIG/Num_Exp/Sedov/ACAT_100.eps}
         \caption{ACAT6 on a $100\times100$ mesh.}
         \label{Sed:AC_100}
     \end{subfigure}
     \hfill
     \begin{subfigure}[b]{0.4\textwidth}
         \centering
         \includegraphics[width=\textwidth]{FIG/Num_Exp/Sedov/ACAT.eps}
         \caption{ACAT6 on a $200\times200$ mesh.}
         \label{Sed:AC}
     \end{subfigure}
     \begin{subfigure}[b]{0.4\textwidth}
         \centering
         \includegraphics[width=\textwidth]{FIG/Num_Exp/Sedov/CATMOODHLLC_100.eps}
         \caption{CATMOOD6 on a $100\times100$ mesh.}
         \label{Sed:CM_100}
     \end{subfigure}
     \hfill
     \begin{subfigure}[b]{0.4\textwidth}
         \centering
         \includegraphics[width=\textwidth]{FIG/Num_Exp/Sedov/CATMOODHLLC.eps}
         \caption{CATMOOD6 on a $200\times200$ mesh.}
         \label{Sed:CM}
     \end{subfigure}
     \caption{Sedov blast wave from section~\ref{ssec:Sedov} --- Scatter plots of the numerical density obtained with a 1st order scheme HLL, and 6th order ACAT6 and CATMOOD6 with HLL flux on the interval $[-1.2,1.2]\times[-1.2,1.2]$.
     Left: results for $100\times100$ uniform cells.
     Right: results for $200\times200$ uniform cells. In black is the exact solution.}
     \label{Se:CM}
\end{figure}
\end{comment}

% --- FIG ---
\begin{figure}[!ht]
     \centering
\begin{tabular}{cccc}
\hspace{-0.5cm}
\rotatebox{90}{\hspace{0.2cm} $100\times 100$ unif. cells} &
\hspace{-0.5cm}
\includegraphics[width=0.35\textwidth]{FIG/Num_Exp/Sedov/HLL_100.eps}&
\hspace{-0.65cm}
\includegraphics[width=0.35\textwidth]{FIG/Num_Exp/Sedov/ACAT_100.eps}&
\hspace{-0.65cm}
\includegraphics[width=0.35\textwidth]{FIG/Num_Exp/Sedov/CATMOODHLLC_100.eps} \\
\hspace{-0.5cm}
\rotatebox{90}{\hspace{0.2cm} $200\times 200$ unif. cells} &
\hspace{-0.5cm}
\includegraphics[width=0.35\textwidth]{FIG/Num_Exp/Sedov/HLL_200.eps}&
\hspace{-0.65cm}
\includegraphics[width=0.35\textwidth]{FIG/Num_Exp/Sedov/ACAT.eps}&
\hspace{-0.65cm}
\includegraphics[width=0.35\textwidth]{FIG/Num_Exp/Sedov/CATMOODHLLC.eps}
\end{tabular}
     \caption{Sedov blast wave from section~\ref{ssec:Sedov} --- Scatter plots of the numerical density obtained with a 1st order scheme HLL, and 6th order ACAT6 and CATMOOD6 with HLL flux on the interval $[-1.2,1.2]\times[-1.2,1.2]$.
     Top: results for $100\times100$ uniform cells.
     Bottom: results for $200\times200$ uniform cells. In black is the exact solution.
     %\raphael{This is a new view of the figure which save space and maybe is more readable.}
     }
     \label{Se:CM}
\end{figure}
% --- FIG ---

\begin{comment}
% --- FIG ---
\begin{figure}[!ht]
     \centering
     \begin{subfigure}[b]{0.495\textwidth}
         \centering
         \includegraphics[width=\textwidth]{FIG/Num_Exp/Sedov/HLLC_100.eps}
         \caption{HLLC on a $100\times100$ mesh.}
         \label{Sed:HLLC_100}
     \end{subfigure}
     \hfill
     \begin{subfigure}[b]{0.495\textwidth}
         \centering
         \includegraphics[width=\textwidth]{FIG/Num_Exp/Sedov/HLLC.eps}
         \caption{HLLC on a $200\times200$ mesh.}
         \label{Sed:HLLC}
     \end{subfigure}
     \caption{Sedov blast wave \ref{ssec:Sedov}. Scatter plots of the numerical solution for the density obtained with HLLC on the interval $[-1.2,1.2]\times[-1.2,1.2]$ adopting a $100\times100$ mesh (left) and $200\times200$ mesh (right).}
     \label{Se:HLLC}
\end{figure}
% --- FIG ---
% --- FIG ---
\begin{figure}[!ht]
     \centering
     \begin{subfigure}[b]{0.495\textwidth}
         \centering
         \includegraphics[width=\textwidth]{FIG/Num_Exp/Sedov/HLL_100.eps}
         \caption{HLL on a $100\times100$ mesh.}
         \label{Sed:HLL_100}
     \end{subfigure}
     \hfill
     \begin{subfigure}[b]{0.495\textwidth}
         \centering
         \includegraphics[width=\textwidth]{FIG/Num_Exp/Sedov/HLL.eps}
         \caption{HLL on a $200\times200$ mesh.}
         \label{Sed:HLL}
     \end{subfigure}
     \caption{Sedov blast wave \ref{ssec:Sedov}. Scatter plots of the numerical solution for the density obtained with HLL on the interval $[-1.2,1.2]\times[-1.2,1.2]$ adopting a $100\times100$ mesh (left) and $200\times200$ mesh (right).}
     \label{Se:HLL}
\end{figure}
% --- FIG ---
% --- FIG ---
\begin{figure}[!ht]
     \centering
     \begin{subfigure}[b]{0.495\textwidth}
         \centering
         \includegraphics[width=\textwidth]{FIG/Num_Exp/Sedov/RUSANOV_100.eps}
         \caption{Rusanov-flux on a $100\times100$ mesh.}
         \label{Sed:Rusanov_100}
     \end{subfigure}
     \hfill
     \begin{subfigure}[b]{0.495\textwidth}
         \centering
         \includegraphics[width=\textwidth]{FIG/Num_Exp/Sedov/RUSANOV.eps}
         \caption{Rusanov-flux on a $200\times200$ mesh.}
         \label{Sed:Rusanov}
     \end{subfigure}
     \caption{Sedov blast wave \ref{ssec:Sedov}. Scatter plots of the numerical solution for the density obtained with Rusanov-flux on the interval $[-1.2,1.2]\times[-1.2,1.2]$ adopting a $100\times100$ mesh (left) and $200\times200$ mesh (right).}
     \label{Se:Rusanov}
\end{figure}
% --- FIG ---
\end{comment}
% --- FIG ---

% Comments
The 2D Sedov problem~\cite{Sedov} is a cylindrical symmetric explosion. The domain is given by $(x,y) \in [-1.2,1.2]^2 $ initially filled with perfect gas at rest such as $(\rho^0, u^0, v^0, p^0, \gamma) = (1, 0, 0, 10^{-13}, 1.4)$. 
%\giovanni{[Why is $p^0$ so low? Where are these numbers taken from? I guess not directly from Sedov's paper.]} 
%\raphael{Indeed when Sedov derived the exact solution he considers $p=0$, that is $\varepsilon=0$. Numerically any 'small' value (viz to $1$) produces the same numerical results in general. }
A total energy of $E_{total}=0.244816$ is concentrated at the origin \cite{LOUBERE2005105}. 
This configuration corresponds to a point-wise symmetric explosion, for which a cylindrical shock front reaches radius $r=\sqrt{x^2+y^2}=1$ at $t_\text{final}=1$ with a density peak of $\rho=6$, see figure~\ref{Sed:CATMOOD6_cut} for an example of a quasi symmetric numerical solution. \\
Figure~\ref{Se:CM} presents the scatter plots of numerical density as a function of cell radius using, respectively for the first and second columns either $100\times100$ or $200\times200$ uniform cells. The tested schemes are  
the 1st order, ACAT6 and CATMOOD6 using HLL numerical flux. 
Notice that the unlimited CAT schemes fails in producing a numerical solution due to the presence of a strong shock wave.
%On CATMOOD6 solution, HLLC scheme has been adopted as first order method. 
%These figures have been obtained identifying the 2D data in a 1D plot O$-rC(r)$ where $r$ is the radius $r\in[0,1.2]$ and $C(r)$ the crowns at distance $r$ and width $dr.$ $C(r) \in \{C_r \in \mathcal{C}\; \rm{such \; that} \;c_r\le C_r < c_{r+dr}\},$ where $\mathcal{C}$ is the family of circumferences. 
% 
Because the exact solution has a cylindrical symmetry (see the black line in figure~\ref{Se:CM}), all cells at the same radius should share the same numerical density, if the scheme exactly preserves the cylindrical symmetry.
The width of the spread/variance of numerical data somehow measures how good the scheme can preserve this symmetry.
As is visible in the figures, the refinement of the grid implies a reduction of the variance for any of the methods.
Moreover, the $6$th order schemes (ACAT6 and CATMOOD6) seem to produce a sharper shock wave and more accurate results. 
Importantly CATMOOD6 has a better cylindrical symmetry than ACAT6.
Of course no positivity issue is reported for any of these schemes.
%Nevertheless, it is possible to observe that only the Rusanov \ref{Sed:Rusanov_100}-\ref{Sed:Rusanov} and HLL \ref{Sed:HLL_100}-\ref{Sed:HLL} method have a small variance. This phenomenon indicates that HLLC \ref{Sed:HLL_100}-\ref{Sed:HLLC} and CAT  methods, both with a posteriori MOOD \ref{Sed:CM_100}-\ref{Sed:CM} technique and a priori ACAT \ref{Sed:AC_100}-\ref{Sed:AC}, are not actually symmetric. Since HLLC was used as first order scheme for both CATMOOD and ACAT methods, to eliminate the possibility that asymmetry is due to HLLC, the same experiment was performed with Rusanov scheme not making any improvements in symmetry. Figure~\ref{Se:CATMOOD6} shows the numerical density obtained with CATMOOD6, where HLLC has been adopted as first order scheme, for the Sedov blast wave problem \ref{ssec:Sedov} on the interval $[-1.2,1.2]\times[-1.2,1.2]$ with $200\times200$ uniform cells at time $t=1$ and CFL$=0.4.$ 
%In particular, figure~\ref{Sed:CATMOOD6} exhibits the projection of the numerical density on the plane O$-xy;$ while figure~\ref{Sed:CATMOOD6_cut} is a zoom if the numerical density on the interval $[0,1.2]\times[0,1.2]$ to show the peak of the numerical solution.



%---------------------------------------------
\subsection{2D Riemann problems} \label{ssec:2DRiemann}
Here, we consider 2D Riemann problems on the interval $[-1,1]\times[-1,1]$ where zero Neumann conditions are applied, and  four initial conditions  are set respectively as in Table~\ref{IC:Riemann} \cite{RP_Schultz93}.
A $400 \times 400$ mesh is adopted for all schemes and configurations. The final time is set to $t_\text{final} = 0.3$ and the CFL to $0.4$.
We run four simulations using CATMOOD6 and the three first order schemes respectively with Rusanov, HLL and HLLC flux functions.
The 1st order scheme with Rusanov flux has been adopted as parachute scheme in the CATMOOD6 cascade.
In these 2D figures the schemes are respectively ordered from top-left to bottom-right.
In addition, in a different set of figures for CATMOOD6 schemes, we plot the percentage of cells updated with CAT6 (top), CAT2 (center) or 1st order Rusanov (bottom) schemes as a function of time-step.
% --- TAB ---
\begin{table}[!ht]
    \centering
    \begin{tabular}{|l l | l l|}
    \hline 
    \multicolumn{4}{|c|}{\textbf{Configuration 3}} \\
    \hline 
   $\rho_2 = 0.5323$  &  $u_2 = 1.206$   &  $\rho_1 = 1.5$    & $u_1 = 0$     \\
   $v_2 = 0$          &  $p_2 = 0.3$     &  $v_1 = 0$         & $p_1 = 1.5$   \\  
    \hline 
    $\rho_3 = 0.138$  &  $u_2 = 1.206$   &  $\rho_4 = 0.5323$ & $u_ 4= 0$     \\
    $v_3 = 1.206$      &  $p_2 = 0.029$   &  $v_4 = 1.206$    & $p_4 = 0.3$  \\  \hline
    \hline 
    \multicolumn{4}{|c|}{\textbf{Configuration 6}} \\
    \hline 
   $\rho_2 = 2     $  &  $u_2 = 0.75 $   &  $\rho_1 = 1.5$    & $u_1 = 0.75$     \\
   $v_2 = 0.5$        &  $p_2 = 1  $     &  $v_1 = -0.5$      & $p_1 = 1$   \\  
    \hline 
   $\rho_3 = 1$       &  $u_2 = -0.75$   &  $\rho_4 = 3$      & $u_4 = -0.75$     \\
   $v_3 = 0.5$        &  $p_2 = 1$       &  $v_4 = -0.5$      & $p_4 =1$  \\  \hline
    \hline 
    \multicolumn{4}{|c|}{\textbf{Configuration 11}} \\
    \hline 
   $\rho_2 = 0.5313$  &  $u_2 = 0.8276$  &  $\rho_1 = 1$      & $u_1 = 0.1$     \\
   $v_2 = 0  $        &  $p_2 = 0.4$     &  $v_1 = 0$         & $p_1 = 1$   \\  
    \hline 
   $\rho_3 = 0.8$     &  $u_2 = 0.1$     &  $\rho_4 = 0.5313$ & $u_4 = 0.1$     \\
   $v_3 = 0$          &  $p_2 = 0.4$     &  $v_4 = 0$         & $p_4 =0.4$  \\  \hline
    \hline 
    \multicolumn{4}{|c|}{\textbf{Configuration 17}} \\
    \hline 
   $\rho_2 = 2     $  &  $u_2 = 0.   $   &  $\rho_1 = 1$      & $u_1 = 0$     \\
   $v_2 = -0.3$       &  $p_2 = 1  $     &  $v_1 = -0.4$      & $p_1 = 1$   \\  
    \hline 
   $\rho_3 = 1.0625$  &  $u_2 = 0$       &  $\rho_4 = 0.5197$ & $u_4 = 0$     \\
   $v_3 = 0.2145$     &  $p_2 = 0.4$     &  $v_4 = -1.1259$   & $p_4 =0.4$  \\  \hline
    \end{tabular} 
    \caption{2D Riemann problem initial conditions. }
    \label{IC:Riemann}
\end{table}
% --- TAB ---

% --- FIG ---
\begin{figure}[!ht]
    \centering
    \begin{subfigure}[b]{0.328\textwidth}
         \centering
         \includegraphics[width=1.3\textwidth]{FIG/Num_Exp/2D_Riemann/Ru_C3_400.eps}
         \caption{Rusanov flux}
         \label{fig:2D_R3_RU}
     \end{subfigure}
     \hfill
     \begin{subfigure}[b]{0.328\textwidth}
    \hspace{-1.cm}
         \centering
         \includegraphics[width=1.3\textwidth]{FIG/Num_Exp/2D_Riemann/HLLC_C3_400.eps}
         \caption{HLLC}
         \label{fig:2D_R3_HLLC}
     \end{subfigure}
     \hfill
     \begin{subfigure}[b]{0.328\textwidth}
    \hspace{-1.cm}
         \centering
         \includegraphics[width=1.3\textwidth]{FIG/Num_Exp/2D_Riemann/CM6_C3_400.eps}
         \caption{CATMOOD6}
         \label{fig:2D_R3_CATMOOD6}
     \end{subfigure}
    \caption{2D Riemann problem from section~\ref{ssec:2DRiemann} with initial condition configuration 3. Zoom of the numerical solution for density on the interval $[-1,1]\times[-1,1]$ adopting a mesh of $400 \times 400-$cells and CFL$=0.4$. Rusanov-flux (a); HLLC (b); and CATMOOD6 with Rusanov flux for the first order method (c).}
   \label{fig:2D_R_C3}
\end{figure}
% --- FIG ---

% --- FIG ---
\begin{figure}[!ht]
    \centering
    \includegraphics[width=0.99\textwidth]{FIG/Num_Exp/2D_Riemann/C3_ord.eps}
    \vspace{-0.3cm}
    \caption{2D Riemann problem from section~\ref{ssec:2DRiemann} with initial condition configuration 3.Percentage of cells updated by CAT6 (top), by CAT2 (center), and by Rusanov (bottom).
    }
   \label{fig:2D_R_C3_ord}
\end{figure}
% --- FIG ---


\begin{comment}
\begin{table}[!ht]
    \centering
    \begin{tabular}{|l l l l|}
    \hline 
    \rm{Lax} & \rm{Configuration 3} &     &   \\
    \hline 
   $\rho_2 = 0.5323$  &  $u_2 = 1.206$   &  $\rho_1 = 1.5$    & $u_1 = 0$     \\
   $v_2 = 0$          &  $p_2 = 0.3$     &  $v_1 = 0$         & $p_1 = 1.5$   \\  
    $\rho_3 = 0.138$  &  $u_2 = 1.206$   &  $\rho_4 = 0.5323$ & $u_ 4= 0$     \\
    $v_3 = 1.206$      &  $p_2 = 0.029$   &  $v_4 = 1.206$    & $p_4 = 0.3$  \\  \hline
    \end{tabular} 
    \caption{ 2D Riemann problem configuration 3. Initial condition. }
    \label{IC:Riemann3}
\end{table}

\begin{table}[!ht]
    \centering
    \begin{tabular}{|l l l l|}
    \hline 
    \rm{Lax} & \rm{Configuration 6} &     &   \\
    \hline 
   $\rho_2 = 2     $  &  $u_2 = 0.75 $   &  $\rho_1 = 1.5$    & $u_1 = 0.75$     \\
   $v_2 = 0.5$        &  $p_2 = 1  $     &  $v_1 = -0.5$      & $p_1 = 1$   \\  
   $\rho_3 = 1$       &  $u_2 = -0.75$   &  $\rho_4 = 3$      & $u_4 = -0.75$     \\
   $v_3 = 0.5$        &  $p_2 = 1$       &  $v_4 = -0.5$      & $p_4 =1$  \\  \hline
    \end{tabular} 
    \caption{ 2D Riemann problem configuration 6. Initial condition. }
    \label{IC:Riemann6}
\end{table}

\begin{table}[!ht]
    \centering
    \begin{tabular}{|l l l l|}
    \hline 
    \rm{Lax} & \rm{Configuration 11} &     &   \\
    \hline 
   $\rho_2 = 0.5313$  &  $u_2 = 0.8276$  &  $\rho_1 = 1$      & $u_1 = 0.1$     \\
   $v_2 = 0  $        &  $p_2 = 0.4$     &  $v_1 = 0$         & $p_1 = 1$   \\  
   $\rho_3 = 0.8$     &  $u_2 = 0.1$     &  $\rho_4 = 0.5313$ & $u_4 = 0.1$     \\
   $v_3 = 0$          &  $p_2 = 0.4$     &  $v_4 = 0$         & $p_4 =0.4$  \\  \hline
    \end{tabular} 
    \caption{ 2D Riemann problem configuration 11. Initial condition. }
    \label{IC:Riemann11}
\end{table}

\begin{table}[!ht]
    \centering
    \begin{tabular}{|l l l l|}
    \hline 
    \rm{Lax} & \rm{Configuration 17} &     &   \\
    \hline 
   $\rho_2 = 2     $  &  $u_2 = 0.   $   &  $\rho_1 = 1$      & $u_1 = 0$     \\
   $v_2 = -0.3$       &  $p_2 = 1  $     &  $v_1 = -0.4$      & $p_1 = 1$   \\  
   $\rho_3 = 1.0625$  &  $u_2 = 0$       &  $\rho_4 = 0.5197$ & $u_4 = 0$     \\
   $v_3 = 0.2145$     &  $p_2 = 0.4$     &  $v_4 = -1.1259$   & $p_4 =0.4$  \\  \hline
    \end{tabular} 
    \caption{ 2D Riemann problem configuration 17. Initial condition. }
    \label{IC:Riemann17}
\end{table}
\end{comment}

\paragraph{Configuration 3}
Figure~\ref{fig:2D_R_C3} shows a zoom of the numerical densities for the configuration 3 (Table~\ref{IC:Riemann}).
All the schemes capture the same global solution. The first order methods are  diffusive even if HLLC scheme performs better. Contrarily, the high order CATMOOD6 clearly improves the sharpness of the shear layers and contacts. 
Figure~\ref{fig:2D_R_C3_ord} exhibits the percentage of cells using CAT6 (top), CAT2 (center) and Rusanov scheme (bottom). We can observe that on average $95\%$ of the cells are updated with $6$th order of accuracy, about $3-4\%$ with $2$nd order or $1$st order. One observes that only few cells demand some limiting, and are sent back to $t^n$ for re-computation by the MOOD approach.
Consequently, the cost of CATMOOD6 is mainly the one of CAT6. 
% --- FIG ---
\begin{figure}[!ht]
    \centering
    \begin{subfigure}[b]{0.328\textwidth}
         \centering
         \includegraphics[width=1.3\textwidth]{FIG/Num_Exp/2D_Riemann/Ru_C6_400.eps}
         \caption{Rusanov flux}
         \label{fig:2D_R6_RU}
     \end{subfigure}
     \hfill
     \begin{subfigure}[b]{0.328\textwidth}
    \hspace{-1.cm}
         \centering
         \includegraphics[width=1.3\textwidth]{FIG/Num_Exp/2D_Riemann/HLLC_C6_400.eps}
         \caption{HLLC}
         \label{fig:2D_R6_HLLC}
     \end{subfigure}
     \hfill
     \begin{subfigure}[b]{0.328\textwidth}
    \hspace{-1.cm}
         \centering
         \includegraphics[width=1.3\textwidth]{FIG/Num_Exp/2D_Riemann/CM6_C6_400.eps}
         \caption{CATMOOD6}
         \label{fig:2D_R6_CATMOOD6}
     \end{subfigure}
    \caption{2D Riemann problem from section~\ref{ssec:2DRiemann} with initial condition configuration 6. Zoom of the numerical solution for density on the interval $[-1,1]\times[-1,1]$ adopting a mesh of $400 \times 400-$cells and CFL$=0.4$. Rusanov-flux (a); HLLC (b); and CATMOOD6 with HLLC for the first order method (c).}
   \label{fig:2D_R_C6}
\end{figure}
% --- FIG ---
% --- FIG ---
\begin{figure}[!ht]
    \centering
    \includegraphics[width=0.99\textwidth]{FIG/Num_Exp/2D_Riemann/C6_ord.eps}
    \vspace{-0.3cm}
    \caption{2D Riemann problem from section~\ref{ssec:2DRiemann} with initial condition configuration 6. Percentage of cells updated by CAT6 (top), by CAT2 (center), and by HLLC (bottom).}
   \label{fig:2D_R_C6_ord}
\end{figure}
% --- FIG ---

\paragraph{Configuration 6}
Likewise for configuration 3, we plot the results in Figure~\ref{fig:2D_R_C6} and \ref{fig:2D_R_C6_ord}.
We again observe that the CATMOOD6 results are far more sharper than the 1st order schemes, with HLLC being the less dissipative scheme of the 1st order ones.
Notice that the shear layers along the curves are Kelvin-Helmholtz unstable, so that it seems normal to see the occurrences of small vortices. This is a numerical evidence of low dissipation of CATMOOD6.
The percentage of troubled cells is of the order of $3-4\%$ on averaged, again showing that CATMOOD6 has globally the same cost than CAT6 plus $10\%$ of the cost of a first order scheme.
% --- FIG ---
\begin{figure}[!ht]
    \centering
    \begin{subfigure}[b]{0.328\textwidth}
         \centering
         \includegraphics[width=1.3\textwidth]{FIG/Num_Exp/2D_Riemann/Ru_C11_400.eps}
         \caption{Rusanov flux}
         \label{fig:2D_R11_RU}
     \end{subfigure}
     \hfill
     \begin{subfigure}[b]{0.328\textwidth}
        \hspace{-1cm}
         \centering
         \includegraphics[width=1.3\textwidth]{FIG/Num_Exp/2D_Riemann/HLLC_C11_400.eps}
         \caption{HLLC}
         \label{fig:2D_R11_HLLC}
     \end{subfigure}
     \hfill
     \begin{subfigure}[b]{0.328\textwidth}
        \hspace{-1cm}
         \centering
         \includegraphics[width=1.3\textwidth]{FIG/Num_Exp/2D_Riemann/CM6_C11_400.eps}
         \caption{CATMOOD6}
         \label{fig:2D_R11_CATMOOD6}
     \end{subfigure}
    \caption{2D Riemann problem from section~\ref{ssec:2DRiemann} with initial condition configuration 11. Zoom of the numerical solution for density on the interval $[-1,1]\times[-1,1]$ adopting a mesh of $400 \times 400-$cells and CFL$=0.4$. Rusanov-flux (a); HLLC (b); and CATMOOD6 with HLLC for the first order method (c).}
   \label{fig:2D_R_C11}
\end{figure}
% --- FIG ---
% --- FIG ---
\begin{figure}[!ht]
    \centering
    \includegraphics[width=0.99\textwidth]{FIG/Num_Exp/2D_Riemann/C11_ord.eps}
    \vspace{-0.3cm}
    \caption{2D Riemann problem from section~\ref{ssec:2DRiemann} with initial condition configuration 11.  Percentage of cells updated by CAT6 (top), by CAT2 (center), and by HLLC (bottom).}
   \label{fig:2D_R_C11_ord}
\end{figure}
% --- FIG ---

\paragraph{Configuration 11}
The numerical solutions for configuration 11 are gathered in Figure~\ref{fig:2D_R_C11} and the percentages of troubled cells in \ref{fig:2D_R_C11_ord}.
It is important to notice that the wave are visibly sharper with the high-order CATMOOD6 as expected. The complex internal pattern is also better captured, avoiding the excessive diffusion that can be seen when any 1st order scheme is employed. 

% --- FIG ---
\begin{figure}[!ht]
    \centering
    \hspace{-0.5cm}
     \begin{subfigure}[b]{0.495\textwidth}
         \centering
        \includegraphics[width=1.0\textwidth]{FIG/Num_Exp/2D_Riemann/Ru_C17_400.eps}
         \caption{Rusanov flux}
         \label{fig:2D_R17_RU}
     \end{subfigure}
     \hfill
     \begin{subfigure}[b]{0.495\textwidth}
         \centering
         \includegraphics[width=1.0\textwidth]{FIG/Num_Exp/2D_Riemann/HLLC_C17_400.eps}
         \caption{HLLC}
         \label{fig:2D_R17_HLLC}
     \end{subfigure}
     \\
    \hspace{-0.5cm}
     \begin{subfigure}[b]{0.495\textwidth}
         \centering
         \includegraphics[width=1.0\textwidth]{FIG/Num_Exp/2D_Riemann/ACAT_C17_400.eps}
         \caption{ACAT6}
         \label{fig:2D_R17_ACAT6}
     \end{subfigure}
     \hfill
     \begin{subfigure}[b]{0.495\textwidth}
         \centering
         \includegraphics[width=1.0\textwidth]{FIG/Num_Exp/2D_Riemann/CM6_C17_400.eps}
         \caption{CATMOOD6}
         \label{fig:2D_R17_CATMOOD6}
     \end{subfigure}
    \caption{2D Riemann problem from section~\ref{ssec:2DRiemann} with initial condition configuration 17. Zoom of the numerical solution for density on the interval $[-1,1]\times[-1,1]$ adopting a mesh of $400 \times 400-$cells and CFL$=0.4$. Rusanov-flux (a); HLLC (b); ACAT6 (c); and CATMOOD6 (d).}
   \label{fig:2D_R_C17}
\end{figure}
% --- FIG ---
% --- FIG ---
\begin{figure}[!ht]
    \centering
    \includegraphics[width=0.99\textwidth]{FIG/Num_Exp/2D_Riemann/C17_ord.eps}
    \vspace{-0.3cm}
    \caption{2D Riemann problem from section~\ref{ssec:2DRiemann} with initial condition configuration 17. Percentage of cells updated by CAT6 (top), by CAT2 (center), and by HLLC (bottom).}
   \label{fig:2D_R_C17_ord}
\end{figure}
% --- FIG ---

\paragraph{Configuration 17}
Figure~\ref{fig:2D_R_C17} and 
Figure~\ref{fig:2D_R_C17_ord} presents the numerical results.
This test presents a small vortex type of structure along with first and secondary waves emanating from the quadruple point.
We can easily observe that CATMOOD6 can capture the vortex and the waves. Contrarily a diffusive scheme would require many more cells to get to this accuracy.
Here the HLLC flux is adopted in the first order scheme of the CATMOOD6 scheme.
The percentage of troubled cells are of the same order than previously, namely $97-98\%$ are updated with $6$th order of accuracy and $2\%$ and $1\%$ with $2$nd or $1$st orders for CATMOOD6. 
The extra cost compared to an unlimited CAT6 scheme is therefore acceptable. 
Table~\ref{CPU:2D_Riemann17} displays, in its first row, the computational cost expressed in seconds for the Configuration 17 concerning Rusanov, HLLC, ACAT6, and CATMOOD6 schemes. Meanwhile, in the second row are presented the ratio between the computational costs with respect to the Rusanov scheme cost. 
Figure~\ref{fig:2D_R_C17} clearly demonstrates the remarkable superiority of the adaptive \textit{a posteriori} approach CAT+MOOD over the \textit{a priori} technique ACAT. 
Despite exhibiting similar computational costs the CATMOOD6 numerical solution significantly outperforms the ACAT6 one in terms of accuracy. 
% --- TAB ----
\begin{table}[!ht]
    \centering
    \numerikNine
    \begin{tabular}{|c||c| c| c| c|}
    \hline 
    & \textbf{Rusanov } & \textbf{HLLC} &  \textbf{ACAT6}   &  \textbf{CATMOOD6} \\
    \hline 
    \textbf{CPU (s)} & 25.28 & 34.11   & 14563.40  & 12871.25   \\ \hline
    \textbf{Ratio} & 1 & 1.35  & 576.08 &  509.15 \\
    \hline
    \end{tabular} 
    \caption{ 2D Riemann problem configuration 17. First row: CPU computational costs expressed in seconds. Second row: ratio between the computational costs with respect to the Rusanov scheme.}
    \label{CPU:2D_Riemann17}
\end{table}
% --- TAB ----


%
% ASTROPHYSICAL JET
%
\clearpage
\subsection{Mach 2000 astrophysical jet}\label{ssec:Mach}
Let us consider the high Mach number astrophysical jet problem \cite{HaShu,ZhangShu2010}. 
% Description
For this test in the interval $[0,1]\times[-0.25,0.25]$ we consider the following initial conditions:
\begin{equation}
\label{IC:Mach2000}
    (\rho^0,u^0,v^0,p^0) = \left\{ \begin{array}{ll}
        (5,800,0,0.4127) & \quad \text{if} \; x=0 \, \text{and} \, y\in [-0.05,0.05], \\ 
        (0.5,0,0,0.4127) & \quad \text{otherwise},
    \end{array} \right.
\end{equation} 
where $\gamma = 5/3$.

\begin{comment}
% --- FIG ---
\begin{figure}[!ht]
    \centering
    \begin{subfigure}[b]{0.495\textwidth}
         \centering
         \includegraphics[width=\textwidth]{FIG/Num_Exp/Mach_2000/M_2000_Ru.eps}
         \caption{Rusanov}
         \label{fig:2D_M2000_Ru}
     \end{subfigure}
     \hfill
     \begin{subfigure}[b]{0.495\textwidth}
         \centering
         \includegraphics[width=\textwidth]{FIG/Num_Exp/Mach_2000/M_2000_HLLC.eps}
         \caption{HLLC}
         \label{fig:2D_M2000_HLLC}
     \end{subfigure}
     \\
     \begin{subfigure}[b]{0.495\textwidth}
         \centering
         \includegraphics[width=\textwidth]{FIG/Num_Exp/Mach_2000/M_2000_CM6.eps}
         \caption{CATMOOD6}
         \label{fig:2D_M2000_CM6}
     \end{subfigure}
    \caption{Mach 2000 from section~\ref{ssec:Mach}. Numerical solution for density on the interval $[0,1]\times[-0.25,0.25]$ adopting a mesh of $300 \times 150-$ quadrangular cells and a CFL$=0.4$. The Rusanov flux (top-left); the HLLC (top-right); and the CATMOOD6 with HLLC for the first order method (bottom).}
   \label{fig:Mach_2000_p}
\end{figure}
% --- FIG ---
% --- FIG ---
\begin{figure}[!ht]
    \centering
    \begin{subfigure}[b]{0.495\textwidth}
         \centering
         \includegraphics[width=\textwidth]{FIG/Num_Exp/Mach_2000/M_2000_Ru_log.eps}
         \caption{Rusanov}
         \label{fig:2D_M2000_Ru_log}
     \end{subfigure}
     \hfill
     \begin{subfigure}[b]{0.495\textwidth}
         \centering
         \includegraphics[width=\textwidth]{FIG/Num_Exp/Mach_2000/M_2000_HLLC_log.eps}
         \caption{HLLC}
         \label{fig:2D_M2000_HLLC_log}
     \end{subfigure}
     \\
     \begin{subfigure}[b]{0.495\textwidth}
         \centering
         \includegraphics[width=\textwidth]{FIG/Num_Exp/Mach_2000/M_2000_CM6_log.eps}
         \caption{CATMOOD6}
         \label{fig:2D_M2000_CM6_log}
     \end{subfigure}
    \caption{Mach 2000 from section~\ref{ssec:Mach}. Numerical solution in the logarithm scale for density on the interval $[0,1]\times[-0.25,0.25]$ adopting a mesh of $300 \times 150-$cells and CFL$=0.4$. The Rusanov flux (top-left); the HLLC (top-right); and the CATMOOD6 with HLLC for the first order method (bottom).}
   \label{fig:Mach_2000_log_p}
\end{figure}
% --- FIG ---
\end{comment}

% --- FIG ---
\begin{figure}[!ht]
    \centering
         \hspace{-0.76cm}
         \includegraphics[width=0.4\textwidth]{FIG/Num_Exp/Mach_2000/M_2000_Ru.eps}
         \hspace{-1.4cm}
         \includegraphics[width=0.4\textwidth]{FIG/Num_Exp/Mach_2000/M_2000_HLLC.eps}
         \hspace{-1.4cm}
         \includegraphics[width=0.4\textwidth]{FIG/Num_Exp/Mach_2000/M_2000_CM6.eps}         
    \caption{Mach 2000 from section~\ref{ssec:Mach}. Numerical solution for density on the interval $[0,1]\times[-0.25,0.25]$ adopting a mesh of $300 \times 150-$ quadrangular cells and a CFL$=0.4$. The Rusanov flux (left); the HLLC (middle); and the CATMOOD6 with HLLC for the first order method (right).}
   \label{fig:Mach_2000_p}
\end{figure}
% --- FIG ---


% --- FIG ---
\begin{figure}[!ht]
    \centering
         \hspace{-0.76cm}
         \includegraphics[width=0.4\textwidth]{FIG/Num_Exp/Mach_2000/M_2000_Ru_log.eps}
         \hspace{-1.4cm}
         \includegraphics[width=0.4\textwidth]{FIG/Num_Exp/Mach_2000/M_2000_HLLC_log.eps}
         \hspace{-1.4cm}
         \includegraphics[width=0.4\textwidth]{FIG/Num_Exp/Mach_2000/M_2000_CM6_log.eps}         
    \caption{Mach 2000 from section~\ref{ssec:Mach}. Numerical solution in the logarithm scale for density on the interval $[0,1]\times[-0.25,0.25]$ adopting a mesh of $300 \times 150-$cells and CFL$=0.4$. The Rusanov flux (Left); the HLLC (middle); and the CATMOOD6 with HLLC for the first order method (right).}
   \label{fig:Mach_2000_log_p}
\end{figure}
% --- FIG ---
% --- FIG ---
\begin{figure}[!ht]
    \centering
    \includegraphics[width=0.99\textwidth]{FIG/Num_Exp/Mach_2000/Mach_2000_ord.eps}
    \vspace{-0.3cm}
    \caption{Mach 2000 from section~\ref{ssec:Mach} with initial condition \eqref{IC:Mach2000}. Percentage of cells updated by CAT6 (top), by CAT2 (center), and by HLLC (bottom).}
   \label{fig:Mach_2000_ord}
\end{figure}
% --- FIG ---
Initially we have a constant ambient gas at rest except for a segment in the left boundary of the domain for which the gas is $10$ times denser and with a high $x$-component of velocity, corresponding to a Mach equals to $2000$. 
Then, we impose inflow boundary conditions in this left segment, and outflow conditions otherwise. 
As such this simulates the penetration of a 
dense jet at hyper-velocity from a portion of the left boundary. This jet generates a bow shock ahead of the jet and a complex shape of the tip of the jet.  Some reference
solutions are reported in \cite{HaShu,ZhangShu2010,BVD20} for instance.
Since the Mach number of the jet is extremely high, negative numerical pressure or density could easily appear during the computation, leading to the crash of the program. 
% 
The simulations are run with CATMOOD6  and some 1st order schemes. The mesh is made of $300 \times 150$ uniform quadrangular cells, and a CFL equals to $0.4$ is adopted for a final time $t_\text{final} = 0.001$. 
% Figures' description
Figures~\ref{fig:Mach_2000_p}-\ref{fig:Mach_2000_log_p} present the numerical densities, with a logarithmic scale in the latter. The simulations are performed with CATMOOD6, first order Rusanov flux and first order HLLC schemes. The first order scheme for the MOOD cascade uses the HLLC numerical flux. 
Notice that the ACAT6 scheme fails for this test due to the generation of nonphysical states (negative pressure). Obviously all unlimited CAP2P schemes for $P\geq1$ also fail.\\
% Comments
The $1$st order scheme can capture the bow shock position but the tip and body of the jet are truly diffused, especially if HLLC flux is not employed. A better shape is gained with HLLC type of flux.
On the contrary CATMOOD6 is able to capture the complexity of the jet motion, its tip and the unstable lateral shear layers producing secondary waves and patterns into the post-shock region. These phenomena would be totally absent from low-accurate simulations for this mesh resolution, hence the need for truly accurate numerical methods. (They are indeed absent from the first order scheme results.)
This phenomenon is even more clear with the logarithmic scale on Figures~\ref{fig:Mach_2000_log_p}.
We would like to emphasize that CATMOOD6 does not have any issue related to positivity, because the parachute scheme is indeed one of the  the $1$st order schemes which is robust enough. As such CAT+MOOD coupling is an almost 'fail-safe' strategy. This is not an obvious property of high-order methods using  classical \apriori limiters. \\
% Percentages
Figure~\ref{fig:Mach_2000_ord} presents the percentage of troubled cells in CATMOOD6 as a function of time-steps.
We plot the cells updated by the schemes in the MOOD cascade, that is with the unlimited CAT6 (top), CAT2 (center) and $1$st order HLLC-base (bottom) schemes.
This is a more advanced unsteady test case compared to the 2D Riemann problems seen in section~\ref{ssec:2DRiemann} for which the solutions were self-similar.
Here, the troubled cell evolution presents two phases: first, for about $500$ time-steps the number of untroubled cells decreases linearly up to $85\%$, and, secondly, this number stagnates for about $250$ time-steps.
Interestingly, the number of cells updated with CAT2 scheme is of the order $10\%$, while truly demanding cells updated with the $1$st order scheme represent about $6\%$. 
These two schemes seem to be useful in our CATMOOD6. \\
% 
This test case is a single example which validates the CATMOOD6 scheme for extremely demanding simulations. Here, both accuracy and robustness are required. 
%In this test case ACAT6 does not provide any solution since it is not robust enough.




